\section{Definición de Escenarios Operativos}
Se definieron tres grupos de prueba (Ramp-up progresivo) atacando simultáneamente múltiples endpoints del sistema:
\begin{table}[h]
\centering
\resizebox{\textwidth}{!}{%
\begin{tabular}{|l|l|l|p{5cm}|l|}
\hline
\textbf{Grupo de Prueba} & \textbf{Nivel de Estrés} & \textbf{Muestras Totales} & \textbf{Endpoints Involucrados} & \textbf{SLA - Hito 1} \\ \hline
\textbf{Grupo 1} & Carga Ligera (Línea Base) & 1,750 pet. (250 x endpoint) & 7 endpoints (Logins, Incidents, Dashboard, Geocoding) & $<$ 3000ms \\ \hline
\textbf{Grupo 2} & Carga Media & 4,500 pet. (750 x endpoint) & 6 endpoints (Logins, Incidents, Dashboard) & $<$ 3000ms \\ \hline
\textbf{Grupo 3} & Carga Pesada (Estrés) & 14,997 pet. ($\sim$2160 x endpoint) & 7 endpoints (Logins, Incidents, Dashboard, Territories) & $<$ 3000ms \\ \hline
\end{tabular}%
}
\caption{Escenarios Operativos de Carga}
\label{tab:escenarios_carga}
\end{table}